\subsection{半单性}

定义 7. --- 设 \(\varrho\) 是拓扑群 G 在希尔伯特空间 E 中的酉表示；称 \(\varrho\) 为半单的，当且仅当存在不可约子表示族 \((\mathrm{F}_i)_{i\in \mathrm{I}}\)，它们是 \(\varrho\) 的子表示，且 E 是子空间 \(\mathrm{F}_i\) 的希尔伯特直和。

有时也说半单酉表示具有离散分解，或称其可离散分解。

如果 \((\varrho_{i})_{i\in\mathrm{I}}\) 是拓扑群 G 的半单酉表示族，那么表示 \(\varrho_{i}\) 的希尔伯特直和是半单的。

命题 7. --- 设 G 是拓扑群，\(\varrho\) 是 G 在复希尔伯特空间 E 中的酉表示。表示 \(\varrho\) 半单，当且仅当 \(\varrho\) 的每个非零子表示都包含不可约子表示。

假设 \(\varrho\) 半单。设 \((\mathrm{F}_i)_{i\in \mathrm{I}}\) 是 \(\varrho\) 的不可约子表示族，使得 E 是子空间 \(\mathrm{F}_i\) 的希尔伯特直和。设 F 是 E 的一个非零子表示。存在 \(i\in \mathrm{I}\)，使得像为 \(\mathrm{F}_i\) 的正交投影在 F 上的限制非零。该正交投影于是通过限制到子空间定义了非零 G-态射 \(p_i\)，从 F 到 \(\mathrm{F}_i\)。由舒尔引理，该 G-态射 \(p_i\) 是满射（V, p. 388 的推论 5）。F 中 \(p_i\) 的核的正交补是 F 的一个与 \(\mathrm{F}_i\) 同构的子表示（V, p. 383 的命题 3），因而不可约。

证明逆命题。令 \(\mathcal{F}\) 为 E 的闭子空间 F 所成的集合，其中 F 在 G 下稳定，且 \(\varrho\) 在 F 中诱导的子表示不可约。令 \(\mathcal{O}\) 为所有满足 \(\mathcal{G}\) 是 \(\mathcal{F}\) 的子集且其成员两两正交的 \(\mathcal{G}\) 所成的集合。

集合 \(\mathcal{O}\) 具有有限特征（E, III, p. 34, 定义 2）。于是令 \(\mathcal{G}\) 为 \(\mathcal{O}\) 的一个极大元（E, III, p. 35, 定理 1）。令 E \(_1\) 为 E 中的子空间，它是 \(\mathcal{G}\) 中不可约子表示 F 的希尔伯特直和。若 E \(_1\) ≠ E，则 E \(_1\) 的正交补非零，而根据假设，G 在 E \(_1^{\circ}\) 中的子表示包含一个不可约子表示。其表示空间 F 与 \(\mathcal{G}\) 中的元素正交，于是 \(\mathcal{G} \cup \{F\} \in \mathcal{O}\)；这与 \(\mathcal{G}\) 的极大性矛盾。因此 E = E \(_1\)，证明完毕。

推论 1. --- 设 G 是拓扑群，\(\varrho\) 是 G 在复希尔伯特空间 E 中的半单酉表示。\(\varrho\) 的每个子表示（以及 \(\varrho\) 的每个商表示）都是半单的。

如果 \(\varrho_{1}\) 是 \(\varrho\) 的子表示，那么 \(\varrho_{1}\) 的每个非零子表示都包含不可约子表示（命题 7），故 \(\varrho_{1}\) 是半单的（同上）。

由 V, p. 383 的命题 3，\(\varrho\) 的每个商表示都与 \(\varrho\) 的一个子表示同构；因此它是半单的。

推论 2. --- 拓扑群 G 的每个有限维酉表示 \(\varrho\) 都是半单的。

这由 V, p. 379 的命题 7 和引理 3 得出。

推论 3. --- 设 \(\varrho_{1}\) 和 \(\varrho_{2}\) 是拓扑群 G 的有限维酉表示。表示 \(\varrho_{1}\) 与 \(\varrho_{2}\) 同构，当且仅当 \(\chi_{\varrho_{1}} = \chi_{\varrho_{2}}\)。

这由推论 2 以及 A, VIII, p. 389 的命题 1, b) 得出。

